% GENERATED FILE. Edit canonical lessons, then run npm run generate:books.
% canonical-source-hash: fnv1a64:a883b3a73017ed99
% canonical-lessons: GU-C305-halvu, GU-C305-dhaglo, GU-C305-dhasvu, GU-C305-datavu, GU-C305-katmal

\chapter{When the Earth Shook}
\label{ch:gu-when-the-earth-shook}
\hlchaptermodality{305}

\begin{chapteropening}
\textbf{By the end of this chapter:} \emph{I can say what an earthquake brought down and who was trapped.}

Complete the last lesson of chapter 305: I can say what an earthquake brought down and who was trapped.
\end{chapteropening}

\section[\texorpdfstring{\gu{હલવું} (halvũ)}{હલવું (halvũ)}]{\gu{હલવું} (halvũ) — to shake; to move}

\label{lesson:GU-C305-halvu}



\begin{quote}
Before the new one: say the Gujarati for a calamity, then the Gujarati for a crisis.
\end{quote}

\subsection*{You'll want to know: \gu{હલવું}}
\textbf{\gu{હલવું}} — \emph{halvũ} — \textquotedblleft{}to shake; to move\textquotedblright{}.

\begin{quote}
\gu{રાત્રે} \gu{અચાનક} \gu{ધરતી} \gu{હલી}.
\end{quote}

\begin{quote}
\emph{rātre acānak dhartī halī.}
\end{quote}

\begin{quote}
Suddenly in the night the ground shook.
\end{quote}

Past (it …): \gu{તે} \gu{મકાન} \gu{હલ્યું}, \emph{te makān halyũ}.

\begin{cousinweb}
\textbf{\gu{હલવું}} (\emph{halvũ}) comes from \emph{*hallati}, \textquotedblleft{}moves, shakes\textquotedblright{}, a form CDIAL reconstructs, where it lists Gujarati \emph{hālvũ}, \textquotedblleft{}to move, shake\textquotedblright{} (CDIAL 14018). \textbf{\gu{હલાવવું}}, to stir, is its causative: you shake a thing, but the ground \textbf{\gu{હલે} \gu{છે}} by itself.
\end{cousinweb}

\begin{grammarlens}[title={What you've built}]
One more word for telling what happened.
\end{grammarlens}

\subsection*{Your turn}
\begin{itemize}
\raggedright
  \item \emph{Say it:} \emph{halvũ}
  \item \emph{Say it:} \emph{halvũ}, once more
  \item \emph{Say it:} \emph{rātre acānak dhartī halī.} — Suddenly in the night the ground shook.
  \item \emph{Recall:} say the Gujarati for a calamity, then the Gujarati for a crisis, then say \emph{halvũ} again
\end{itemize}

\begin{tcolorbox}[breakable,colback=teal!4,colframe=teal!35!black,title={Before you move on}]
What does \gu{હલવું} mean? (\textquotedblleft{}To shake; to move\textquotedblright{}.) Say it once more.
\end{tcolorbox}
\section[\texorpdfstring{\gu{ઢગલો} (ḍhaglo)}{ઢગલો (ḍhaglo)}]{\gu{ઢગલો} (ḍhaglo) — a heap}

\label{lesson:GU-C305-dhaglo}



\begin{quote}
Before the new one: say the Gujarati for a crisis, then the Gujarati for to shake.
\end{quote}

\subsection*{You'll want to know: \gu{ઢગલો}}
\textbf{\gu{ઢગલો}} — \emph{ḍhaglo} — \textquotedblleft{}a heap\textquotedblright{}.

\begin{quote}
\gu{ભૂકંપ} \gu{પછી} \gu{રસ્તા} \gu{પર} \gu{પથ્થરનો} \gu{ઢગલો} \gu{હતો}.
\end{quote}

\begin{quote}
\emph{bhūkamp pachī rastā par paththarno ḍhaglo hato.}
\end{quote}

\begin{quote}
After the earthquake there was a heap of stones on the road.
\end{quote}

A masculine noun.

\begin{cousinweb}
\textbf{\gu{ઢગલો}} (\emph{ḍhaglo}) comes from \emph{*ḍhikka}, a form CDIAL reconstructs, where it lists Gujarati \emph{ḍhaglo}, \textquotedblleft{}a heap\textquotedblright{} (CDIAL 5585). It is masculine: \textbf{\gu{પથ્થરનો} \gu{ઢગલો}}, a heap of stones.
\end{cousinweb}

\begin{grammarlens}[title={What you've built}]
One more word for telling what happened.
\end{grammarlens}

\subsection*{Your turn}
\begin{itemize}
\raggedright
  \item \emph{Say it:} \emph{ḍhaglo}
  \item \emph{Say it:} \emph{ḍhaglo}, once more
  \item \emph{Say it:} \emph{bhūkamp pachī rastā par paththarno ḍhaglo hato.} — After the earthquake there was a heap of stones on the road.
  \item \emph{Recall:} say the Gujarati for a crisis, then the Gujarati for to shake, then say \emph{ḍhaglo} again
\end{itemize}

\begin{tcolorbox}[breakable,colback=teal!4,colframe=teal!35!black,title={Before you move on}]
What does \gu{ઢગલો} mean? (\textquotedblleft{}A heap\textquotedblright{}.) Say it once more.
\end{tcolorbox}
\section[\texorpdfstring{\gu{ધસવું} (dhasvũ)}{ધસવું (dhasvũ)}]{\gu{ધસવું} (dhasvũ) — to give way; to cave in}

\label{lesson:GU-C305-dhasvu}



\begin{quote}
Before the new one: say the Gujarati for to shake, then the Gujarati for a heap.
\end{quote}

\subsection*{You'll want to know: \gu{ધસવું}}
\textbf{\gu{ધસવું}} — \emph{dhasvũ} — \textquotedblleft{}to give way; to cave in\textquotedblright{}.

\begin{quote}
\gu{ભૂકંપમાં} \gu{જૂનું} \gu{મકાન} \gu{ધસી} \gu{પડ્યું}.
\end{quote}

\begin{quote}
\emph{bhūkampmā̃ jūnũ makān dhasī paḍyũ.}
\end{quote}

\begin{quote}
The old building caved in during the earthquake.
\end{quote}

Past (it …): \gu{તે} \gu{મકાન} \gu{ધસી} \gu{પડ્યું}, \emph{te makān dhasī paḍyũ}.

\begin{cousinweb}
\textbf{\gu{ધસવું}} (\emph{dhasvũ}) is inherited from Sanskrit \emph{dhváṁsati}, \textquotedblleft{}falls to pieces\textquotedblright{}, and CDIAL lists the Gujarati verb, \textquotedblleft{}to push through, be undermined\textquotedblright{} (CDIAL 6896). With \textbf{\gu{પડવું}} it is a building giving way: \textbf{\gu{ધસી} \gu{પડ્યું}}.
\end{cousinweb}

\begin{grammarlens}[title={What you've built}]
One more word for telling what happened.
\end{grammarlens}

\subsection*{Your turn}
\begin{itemize}
\raggedright
  \item \emph{Say it:} \emph{dhasvũ}
  \item \emph{Say it:} \emph{dhasvũ}, once more
  \item \emph{Say it:} \emph{bhūkampmā̃ jūnũ makān dhasī paḍyũ.} — The old building caved in during the earthquake.
  \item \emph{Recall:} say the Gujarati for to shake, then the Gujarati for a heap, then say \emph{dhasvũ} again
\end{itemize}

\begin{tcolorbox}[breakable,colback=teal!4,colframe=teal!35!black,title={Before you move on}]
What does \gu{ધસવું} mean? (\textquotedblleft{}To give way; to cave in\textquotedblright{}.) Say it once more.
\end{tcolorbox}
\section[\texorpdfstring{\gu{દટાવું} (daṭāvũ)}{દટાવું (daṭāvũ)}]{\gu{દટાવું} (daṭāvũ) — to be buried (under something)}

\label{lesson:GU-C305-datavu}



\begin{quote}
Before the new one: say the Gujarati for a heap, then the Gujarati for to give way.
\end{quote}

\subsection*{You'll want to know: \gu{દટાવું}}
\textbf{\gu{દટાવું}} — \emph{daṭāvũ} — \textquotedblleft{}to be buried (under something)\textquotedblright{}.

\begin{quote}
\gu{ભૂકંપમાં} \gu{ઘણા} \gu{લોકો} \gu{મકાનો} \gu{નીચે} \gu{દટાઈ} \gu{ગયા}.
\end{quote}

\begin{quote}
\emph{bhūkampmā̃ ghaṇā loko makāno nīce daṭāī gayā.}
\end{quote}

\begin{quote}
In the earthquake many people were buried under buildings.
\end{quote}

Past (he …): \gu{તે} \gu{માણસ} \gu{દટાઈ} \gu{ગયો}, \emph{te māṇas daṭāī gayo}.

\begin{cousinweb}
\textbf{\gu{દટાવું}} (\emph{daṭāvũ}) is the intransitive partner of \emph{dāṭvũ}, to bury, which CDIAL lists under \emph{*draṭ}, \textquotedblleft{}press\textquotedblright{}, a form it reconstructs (CDIAL 6618): to be pressed down out of sight. People and things \textbf{\gu{દટાઈ} \gu{જાય} \gu{છે}} under earth or fallen walls.
\end{cousinweb}

\begin{grammarlens}[title={What you've built}]
One more word for telling what happened.
\end{grammarlens}

\subsection*{Your turn}
\begin{itemize}
\raggedright
  \item \emph{Say it:} \emph{daṭāvũ}
  \item \emph{Say it:} \emph{daṭāvũ}, once more
  \item \emph{Say it:} \emph{bhūkampmā̃ ghaṇā loko makāno nīce daṭāī gayā.} — In the earthquake many people were buried under buildings.
  \item \emph{Recall:} say the Gujarati for a heap, then the Gujarati for to give way, then say \emph{daṭāvũ} again
\end{itemize}

\begin{tcolorbox}[breakable,colback=teal!4,colframe=teal!35!black,title={Before you move on}]
What does \gu{દટાવું} mean? (\textquotedblleft{}To be buried (under something)\textquotedblright{}.) Say it once more.
\end{tcolorbox}
\section[\texorpdfstring{\gu{કાટમાળ} (kāṭmāḷ)}{કાટમાળ (kāṭmāḷ)}]{\gu{કાટમાળ} (kāṭmāḷ) — rubble}

\label{lesson:GU-C305-katmal}



\begin{quote}
Before the new one: say the Gujarati for to give way, then the Gujarati for to be buried.
\end{quote}

\subsection*{You'll want to know: \gu{કાટમાળ}}
\textbf{\gu{કાટમાળ}} — \emph{kāṭmāḷ} — \textquotedblleft{}rubble\textquotedblright{}.

\begin{quote}
\gu{લોકોએ} \gu{કાટમાળ} \gu{નીચેથી} \gu{એક} \gu{બાળકને} \gu{કાઢ્યું}.
\end{quote}

\begin{quote}
\emph{lokoe kāṭmāḷ nīcethī ek bāḷakne kāḍhyũ.}
\end{quote}

\begin{quote}
People pulled a child out from under the rubble.
\end{quote}

A masculine noun.

\begin{cousinweb}
\textbf{\gu{કાટમાળ}} (\emph{kāṭmāḷ}) No dictionary we use records an older form, so its history is not settled. It is what is left of a building that has fallen: broken walls, stones and wood.
\end{cousinweb}

\begin{grammarlens}[title={What you've built}]
One more word for telling what happened.
\end{grammarlens}

\subsection*{Your turn}
\begin{itemize}
\raggedright
  \item \emph{Say it:} \emph{kāṭmāḷ}
  \item \emph{Say it:} \emph{kāṭmāḷ}, once more
  \item \emph{Say it:} \emph{lokoe kāṭmāḷ nīcethī ek bāḷakne kāḍhyũ.} — People pulled a child out from under the rubble.
  \item \emph{Recall:} say the Gujarati for to give way, then the Gujarati for to be buried, then say \emph{kāṭmāḷ} again
\end{itemize}

\begin{tcolorbox}[breakable,colback=teal!4,colframe=teal!35!black,title={Before you move on}]
What does \gu{કાટમાળ} mean? (\textquotedblleft{}Rubble\textquotedblright{}.) Say it once more.
\end{tcolorbox}
